% =====================================================================
%  Artigo de mestrado — Biomassa acima do solo (AGB) na Amazônia por LiDAR
%  Formatado em conformidade com a ABNT (classe abntex2, modo "article" —
%  NBR 6022 artigo científico; NBR 10520 citações autor-data; NBR 6023
%  referências; NBR 14724 margens/espaçamento/tipografia).
%
%  ESTRUTURA DO PROJETO
%  ---------------------
%  main.tex            — este arquivo: preâmbulo + pré-textual + \input das
%                         seções + referências.
%  secoes/01-...tex     — uma seção do artigo por arquivo.
%  referencias.tex      — lista de referências (NBR 6023), em arquivo próprio.
%  figuras/              — imagens do artigo (ainda vazia; ver figuras/LEIA-ME.txt).
%
%  COMO COMPILAR
%  --------------
%  tectonic -X compile main.tex   (ou VS Code + LaTeX Workshop, autobuild já
%  configurado em .vscode/settings.json, mas abra a PASTA do projeto, não só
%  o arquivo, senão o autobuild não liga).
%
%  As Figuras 1–4 e a Tabela 1 (comparação de modelos) estão como
%  PLACEHOLDERS — procure por "PLACEHOLDER" dentro de secoes/ e troque pelo
%  \includegraphics ou pela tabela final quando exportar os resultados.
% =====================================================================
\documentclass[article,12pt,oneside,a4paper]{abntex2}

\usepackage[utf8]{inputenc}
\usepackage[T1]{fontenc}
\usepackage{lmodern}
\usepackage{amsmath,amssymb}
\usepackage{graphicx}
\usepackage{array}
\usepackage{booktabs}
\usepackage{caption}
\usepackage{float}
\usepackage{csquotes}
\usepackage[normalem]{ulem}
\usepackage{xurl}

\graphicspath{{figuras/}}

% Caixa de placeholder para figuras/tabelas ainda não inseridas
\newcommand{\figplaceholder}[1]{%
  \fbox{\parbox[c][4cm][c]{0.8\textwidth}{\centering\itshape #1}}%
}

\titulo{Estrutura vertical do dossel e estimativa de biomassa acima do solo na
  Amazônia por LiDAR aerotransportado: parcimônia versus profundidade sob
  amostra reduzida}
\tituloestrangeiro{Canopy vertical structure and aboveground biomass
  estimation in the Amazon by airborne LiDAR: parsimony versus depth under a
  reduced sample}
\autor{Daniel Jacob}
\local{Brasil}
\data{2026}

\begin{document}

\maketitle

\begin{resumoumacoluna}
\noindent
\textbf{Contexto.} A biomassa acima do solo (AGB) é a variável de referência
para quantificar estoques de carbono florestal, mas sensores ópticos orbitais
saturam em florestas tropicais densas por enxergarem apenas o topo do dossel.
O LiDAR aerotransportado mede diretamente a estrutura vertical da vegetação ---
a dimensão com maior poder preditivo de biomassa --- e permite calibrar relações
entre estrutura e AGB de campo que, em tese, se estendem a áreas sem inventário.

\textbf{Objetivo.} Este trabalho pergunta, especificamente: dentro de um
conjunto de parcelas de inventário florestal amazônico com cobertura LiDAR
pareada, qual fração da variância na AGB medida em campo é explicada pela
estrutura do dossel --- e uma representação parcimoniosa dessa estrutura
(estatísticas agregadas de altura) generaliza tão bem quanto arquiteturas de
aprendizado profundo aplicadas à geometria bruta da nuvem de pontos?

\textbf{Dados e método.} Cruzaram-se nuvens de pontos LiDAR do programa NASA CMS
Brazil (ORNL DAAC) com inventários florestais de campo em 28 sites amazônicos,
resultando em \textbf{493 parcelas pareáveis} (AGB de campo $\times$ nuvem de
pontos LiDAR) após auditoria e correção de bugs de geometria e pareamento ---
\textbf{325} após remoção de outliers de área e forma. A biomassa de referência
foi calculada árvore a árvore por equações alométricas (Chave et al. 2014;
Chambers et al. 2000; Goodman et al. 2013) e uma correção do descompasso
temporal entre voo LiDAR e medição de campo. Seis abordagens de modelagem foram
comparadas sob validação cruzada k-fold: um GBR sobre a distribuição completa de
alturas do dossel (1025 atributos), sua variante com redução de dimensionalidade
(GBR sobre $K$ quantis, $K \in \{8,\dots,256\}$), um GBR e um Random Forest sobre
37 métricas estruturais clássicas (\textit{Area-Based Approach}), uma PointNet
sobre a nuvem crua, e duas CNNs --- 2D sobre rasters de altura e 3D sobre voxels
de ocupação.

\textbf{Resultados.} Testando cada tratamento de pré-processamento isoladamente
contra a mesma nuvem crua --- projeção do solo, filtro de altura $\leq$1,3~m,
remoção de outliers geométricos de parcela, e correção de gap temporal, cada um
sozinho, sem empilhar com os demais (Seção~\ref{sec:preproc}) --- a remoção de
outliers isolada se mostrou, disparado, o tratamento que mais importa: sozinha,
já supera o pipeline com todos os tratamentos combinados em quatro dos sete
modelos. Projeção do solo e filtro de altura, isolados, têm efeito pequeno; a
correção de gap temporal isolada é, na maioria dos modelos, mais fraca que não
aplicar tratamento nenhum --- ela só ``ajuda'' quando empilhada em cima da
remoção de outliers, cujo efeito domina. Na melhor variante de cada modelo ---
que para cinco deles é justamente a remoção de outliers isolada ---, a CNN 2D
sobre rasters de altura obteve o melhor resultado ($R^2$ = 0,547;
rRMSE $\approx$ 41,7\%), seguida por Random Forest ($R^2$ = 0,532), GBR sobre a
distribuição completa de alturas ($R^2$ = 0,526), ABA ($R^2$ = 0,514) e GBR
sobre quantis de altura ($R^2$ = 0,510); a CNN 3D sobre voxels (que só atinge seu
melhor resultado com o pipeline combinado completo, não isolado) ficou abaixo de
todas essas ($R^2$ = 0,455); a PointNet sobre a nuvem crua (idem, precisa do
pipeline combinado) foi consistentemente a mais fraca ($R^2 \approx$ 0,21). Já na
ablação de dimensionalidade, \textbf{8 quantis de altura bastam para superar o
vetor completo de 1024 alturas} ($R^2$ = 0,43 vs. 0,39--0,42), e o ganho
adicional até 256 quantis é marginal --- evidência de que a maior parte da alta
dimensão original era ruído alimentando sobreajuste, não sinal.

\textbf{Conclusão.} A estrutura vertical do dossel, resumida de forma simples,
explica uma fração moderada e competitiva (com a literatura de LiDAR aéreo
tropical e com o produto orbital GEDI) da variância em AGB --- e fá-lo com
modelos ordens de grandeza mais simples que arquiteturas profundas sobre
geometria bruta. Este resultado é lido com uma ressalva central discutida no
corpo do artigo: a validação cruzada aqui empregada é k-fold aleatório, não
\textit{leave-one-site-out}, e portanto mede ajuste dentro da amostra de sítios
já vistos --- não generalização para floresta amazônica não amostrada, que é o
padrão mais alto exigido por um uso operacional em contabilidade de carbono.

\vspace{\onelineskip}
\noindent\textbf{Palavras-chave}: biomassa acima do solo; LiDAR aerotransportado;
Amazônia; aprendizado de máquina; redução de dimensionalidade; GEDI.
\end{resumoumacoluna}

\begin{resumoumacoluna}[Abstract]
\begin{otherlanguage*}{english}
\noindent
\textbf{Context.} Aboveground biomass (AGB) is the reference variable for
quantifying forest carbon stocks, but optical orbital sensors saturate in
dense tropical forests because they only see the top of the canopy. Airborne
LiDAR directly measures the vertical structure of vegetation --- the
dimension with the greatest predictive power for biomass --- and allows the
calibration of structure--AGB relationships from field data that, in
principle, extend to areas without inventory.

\textbf{Objective.} This work asks, specifically: within a set of Amazonian
forest-inventory plots with paired LiDAR coverage, what fraction of the
variance in field-measured AGB is explained by canopy structure --- and does a
parsimonious representation of that structure (aggregated height statistics)
generalize as well as deep-learning architectures applied to raw point-cloud
geometry?

\textbf{Data and methods.} LiDAR point clouds from the NASA CMS Brazil program
(ORNL DAAC) were crossed with field forest inventories across 28 Amazonian
sites, yielding \textbf{493 pairable plots} (field AGB $\times$ LiDAR point
cloud) after auditing and fixing geometry/pairing bugs --- \textbf{325} after
removing area/shape outliers. Reference biomass was computed tree-by-tree
using allometric equations (Chave et al. 2014; Chambers et al. 2000; Goodman
et al. 2013) and a correction for the temporal mismatch between the LiDAR
flight and the field measurement. Six modeling approaches were compared under
k-fold cross-validation: a GBR over the full canopy-height distribution (1,025
features), its dimensionality-reduced variant (GBR over $K$ height quantiles,
$K \in \{8,\dots,256\}$), a GBR and a Random Forest over 37 classic structural
metrics (\textit{Area-Based Approach}), a PointNet over the raw point cloud,
and two CNNs --- 2D over height rasters and 3D over occupancy voxels.

\textbf{Results.} Testing each preprocessing treatment in isolation against
the same raw cloud --- ground projection, height filter $\leq$1.3~m, removal
of geometric plot outliers, and temporal-gap correction, each alone, without
stacking the others --- outlier removal alone proved by far the treatment
that matters most: alone, it already surpasses the pipeline with all
treatments combined in four of seven models. Ground projection and the height
filter, in isolation, have a small effect; isolated temporal-gap correction
is, in most models, weaker than applying no treatment at all --- it only
``helps'' when stacked on top of outlier removal, whose effect dominates. In
the best variant of each model --- which for five of them is precisely
isolated outlier removal --- the 2D CNN over height rasters obtained the best
result ($R^2$ = 0.547; rRMSE $\approx$ 41.7\%), followed by Random Forest
($R^2$ = 0.532), GBR over the full height distribution ($R^2$ = 0.526), ABA
($R^2$ = 0.514) and GBR over height quantiles ($R^2$ = 0.510); the 3D CNN over
voxels (which only reaches its best result with the full combined pipeline,
not isolated) fell below all of these ($R^2$ = 0.455); the PointNet over the
raw cloud (likewise needing the combined pipeline) was consistently the
weakest ($R^2 \approx$ 0.21). In the dimensionality ablation, \textbf{8 height
quantiles are enough to beat the full vector of 1,024 heights} ($R^2$ = 0.43
vs. 0.39--0.42), and the additional gain up to 256 quantiles is marginal ---
evidence that most of the original high dimensionality was noise feeding
overfitting, not signal.

\textbf{Conclusion.} Canopy vertical structure, summarized simply, explains a
moderate and competitive fraction (relative to the tropical airborne-LiDAR
literature and to the GEDI orbital product) of the variance in AGB --- and
does so with models orders of magnitude simpler than deep architectures over
raw geometry. This result is read with a central caveat discussed in the body
of the article: the cross-validation used here is random k-fold, not
leave-one-site-out, and therefore measures within-sample fit to already-seen
sites --- not spatial generalization to unsampled Amazonian forest, which is
the higher bar required for operational use in carbon accounting.

\vspace{\onelineskip}
\noindent\textbf{Keywords}: aboveground biomass; airborne LiDAR; Amazon;
machine learning; dimensionality reduction; GEDI.
\end{otherlanguage*}
\end{resumoumacoluna}

% !TEX root = ../main.tex
\section{Introdução}
\label{sec:intro}

Cerca de metade da biomassa vegetal é carbono, fração de referência adotada
pelo \textit{Intergovernmental Panel on Climate Change} para inventários
nacionais de gases de efeito estufa (IPCC, 2019). As florestas tropicais
intactas armazenam, globalmente, parte substancial desse estoque terrestre, e
a Amazônia concentra uma das maiores reservas de carbono florestal do planeta
(PAN et al., 2011; SAATCHI et al., 2011; FAO, 2020). Quantificar a biomassa
acima do solo (\textit{aboveground biomass}, AGB) com precisão espacial é o
que sustenta inventários nacionais de emissões segundo as diretrizes
metodológicas do IPCC (IPCC, 2019), mecanismos de compensação como o REDD+,
mercados de crédito de carbono e o acompanhamento de metas do Acordo de Paris
(UNFCCC, 2015). Sem estimativas confiáveis e espacialmente distribuídas, não é
possível distinguir floresta intacta de degradada, dimensionar perdas por
desmatamento, nem validar os modelos do ciclo global de carbono que dependem
desses números como entrada (IPCC, 2021).

O gargalo histórico dessa estimativa é que AGB depende fortemente da
\textbf{estrutura vertical} da floresta --- altura das árvores, densidade e
arquitetura do dossel ---, exatamente a dimensão que sensores ópticos orbitais
(Landsat, Sentinel) enxergam mal: eles registram a reflectância do topo do
dossel e saturam em florestas densas de alta biomassa, como a maior parte da
Amazônia intacta. O LiDAR (\textit{Light Detection and Ranging}) resolve esse
problema na origem --- pulsos laser penetram o dossel e retornam um perfil
tridimensional da vegetação, medindo diretamente altura e estrutura vertical.
Isso o torna substancialmente mais acurado que índices espectrais em floresta
tropical de alta biomassa e permite calibrar relações alométricas (métricas
LiDAR $\leftrightarrow$ AGB de campo) que, em princípio, podem ser estendidas no
espaço e no tempo --- reduzindo a dependência de inventários de campo, caros e
logisticamente inviáveis na escala da bacia amazônica.

Para operacionalizar esse objetivo, a biomassa de referência desta tese é
calculada pelas mesmas equações alométricas que sustentam boa parte da
literatura de carbono florestal tropical: Chave et al. (2014) para árvores
vivas, Chambers et al. (2000) para árvores mortas em pé e Goodman et al.
(2013) para palmeiras --- a mesma base alométrica usada, por exemplo, por
Longo et al. (2016) na Amazônia.

A literatura de estimativa de AGB por LiDAR aerotransportado em floresta tropical
é madura e compartilha um padrão metodológico que vale destacar frente a este
trabalho: Asner \& Mascaro (2014) --- 804 parcelas de 1~ha em quatro países
tropicais, $R^2$ = 0,90 usando altura de topo de dossel (TCH) combinada a área
basal e densidade de madeira --- validam por \textbf{sítio}, não por parcela
embaralhada aleatoriamente. Longo et al. (2016), na Amazônia brasileira (359
parcelas, intacta e degradada), reporta $R^2 \approx$ 0,70 por regressão sobre
métricas LiDAR agregadas --- mais próximo em escopo e escala deste trabalho, mas
também sem validação k-fold aleatória explícita. Bispo et al. (2020) --- atribuída
ao Cerrado, não à Amazônia, após a correção de erro desta auditoria
(Seção~\ref{sec:lit-check}) --- reporta RMSE = 67,6~Mg/ha, rRMSE = 36,0\%,
$R^2$ = 0,58 usando só métricas LiDAR, números na mesma ordem de grandeza dos
obtidos aqui (RMSE $\approx$ 73--79~Mg/ha, rRMSE $\approx$ 43--46\%, variante sem
outliers), apesar do bioma diferente. d'Oliveira et al. (2012), em floresta
manejada no Acre, reporta $R^2$ = 0,63--0,72 em escala de \textbf{talhão}
(1.000~ha) --- suporte espacial muito maior que a parcela individual usada aqui,
o que limita a comparabilidade direta apesar da mesma região amazônica.

Mais recentemente, a literatura passou a incorporar aprendizado profundo
diretamente sobre a nuvem de pontos bruta: Oehmcke et al. (2024), em floresta
temperada dinamarquesa, encontra uma CNN de Minkowski superando PointNet, KPConv
e uma abordagem \textit{area-based} clássica --- o oposto do padrão encontrado
neste trabalho (Seção~\ref{sec:arch}), em que a abordagem parcimoniosa vence a
PointNet por uma margem grande. A leitura mais provável não é uma contradição
entre os dois resultados, mas um efeito de \textbf{regime de amostra}: florestas
temperadas manejadas tendem a ter conjuntos de treino bem maiores e menos
heterogêneos estruturalmente que os 325 pares LiDAR-inventário amazônicos
disponíveis aqui --- consistente com a leitura central deste trabalho: sob
orçamento de amostra reduzido, o efeito da complexidade representacional da
arquitetura compete com --- e pode ser dominado por --- decisões de
pré-processamento e engenharia de \textit{feature} do dado de entrada.

Os produtos orbitais da NASA (GEDI L4A/L4B) operam em escala e sensor distintos
--- LiDAR espacial \textit{footprint} ($\sim$25~m) ou grade de 1~km, não nuvem
aérea densa por parcela --- e servem aqui como baliza de ordem de grandeza, não
como comparação direta: a validação independente do GEDI L4A contra LiDAR aéreo
em savana africana (Li et al., 2024) reporta $R^2 \approx$ 0,42, abaixo de todos
os modelos desta tese na variante sem outliers, mas sobre um problema de
estimação substancialmente mais difícil (resolução espacial muito menor, amostra
espacialmente esparsa).

Um ponto, porém, permanece pouco caracterizado empiricamente para o contexto
amazônico. A literatura revisada acima mistura, na prática, dois eixos que
raramente são testados em separado: \textbf{de um lado}, o quanto a arquitetura
de modelagem em si --- de estatísticas agregadas simples a redes profundas
sobre geometria bruta --- importa para o desempenho; \textbf{de outro}, o
quanto decisões de pré-processamento e engenharia de \textit{feature} do dado
de entrada (curadoria geométrica de parcela, projeção de relevo, filtros de
altura, correção temporal) importam. Quando os dois efeitos são empilhados num
único pipeline, como é comum na literatura, não é possível atribuir o ganho de
desempenho observado a um ou a outro --- e, sob um orçamento de amostra
pequeno (nesta tese, no máximo 493 parcelas, 325 após controle de qualidade
geométrica), essa indistinção importa: um resultado que parece confirmar o
valor de uma arquitetura sofisticada pode, na verdade, ser inteiramente
explicado por uma etapa de curadoria de dados anterior ao modelo.

\label{sec:rq}
Este trabalho investiga uma pergunta específica, deliberadamente mais restrita do
que a aspiração de ``mapear a biomassa da Amazônia'':

\begin{quote}
\textbf{Num regime de poucas centenas de parcelas rotuladas, o que domina a
capacidade de um modelo generalizar na estimativa de AGB por LiDAR --- a
complexidade representacional da arquitetura ou o estágio de
pré-processamento/engenharia de \textit{feature} do dado de entrada --- e dá
para separar esses dois efeitos?}
\end{quote}

A hipótese de trabalho, motivada pela teoria estatística clássica de
viés-variância sob alta dimensionalidade e poucas amostras, é que, neste
regime de amostra reduzida, o estágio de curadoria e pré-processamento do dado
de entrada tem efeito de magnitude igual ou maior que a escolha da
arquitetura de modelagem --- e que representações de alta capacidade (PointNet
sobre nuvem crua, CNN 3D sobre voxels) sofrem desproporcionalmente mais de
sobreajuste que representações parcimoniosas (estatísticas agregadas de
altura), independentemente da riqueza estrutural que a geometria bruta, em
princípio, carrega. Testar essa hipótese exige um desenho experimental que
isole os dois eixos --- pré-processamento e arquitetura --- em vez de tratá-los
como uma correção única, como detalhado no protocolo de validação
(Seção~\ref{sec:protocolo}).

Essa pergunta se desdobra em quatro objetivos específicos, cada um respondido
numa subseção dedicada dos Resultados (Seção~\ref{sec:resultados}):

\begin{enumerate}
  \item \textbf{Eixo pré-processamento}: isolar o efeito do controle de
    qualidade geométrica das parcelas e da correção do gap temporal entre voo
    LiDAR e medição de campo sobre o desempenho dos modelos, testando cada
    tratamento isoladamente em vez de empilhá-los numa correção única
    (\textbf{§\ref{sec:preproc}}).
  \item \textbf{Eixo arquitetura}: comparar, sob o mesmo protocolo de
    validação e a mesma variante de pré-processamento, seis abordagens de
    modelagem que vão de estatísticas agregadas interpretáveis (ABA, Random
    Forest) a aprendizado profundo sobre a geometria bruta (PointNet, CNN 2D,
    CNN 3D), identificando qual representação generaliza melhor neste regime de
    amostra pequena (\textbf{§\ref{sec:arch}}).
  \item \textbf{Dentro do eixo arquitetura}, em grão mais fino: medir como o
    desempenho do GBR varia com a dimensionalidade da representação de altura
    do dossel (de 8 a 1024 atributos), para testar se mais informação bruta
    realmente se traduz em mais sinal preditivo ou se é absorvida como ruído
    de sobreajuste (\textbf{§\ref{sec:ablation}}).
  \item Identificar, através do único modelo desta comparação com importância de
    variável interpretável (ABA), quais propriedades estruturais do dossel mais
    sustentam a previsão de AGB (\textbf{§\ref{sec:interpret}}).
\end{enumerate}

Por rigor, delimita-se desde já explicitamente o que \textbf{não} está dentro do escopo do
resultado central: a validação aqui é k-fold aleatório, não
\textit{leave-one-site-out} --- portanto mede ajuste dentro da amostra de sítios
já vistos no treino, não generalização espacial para floresta amazônica não
amostrada. Há, agora, um intervalo de predição por parcela (conformal prediction
sobre resíduos out-of-fold, Seção~\ref{sec:conformal}) --- mas ele cobre só a
incerteza do modelo estrutura$\rightarrow$AGB, não o erro alométrico da fórmula de
referência nem a incerteza de medição de campo propagados até o alvo de treino.
Uma análise de sensibilidade (Seção~\ref{sec:allom-sens}) mostra, além disso, que
o próprio resultado ``CNN 2D é a melhor arquitetura'' não é robusto à escolha da
fórmula alométrica de referência --- sob densidade de madeira específica por
espécie (M3) em vez de uniforme (M1), o Random Forest assume a liderança. Essas
lacunas são discutidas com profundidade na Seção~\ref{sec:discussao} e não são
triviais: são precisamente o que separa ``a estrutura do dossel explica biomassa
nesta amostra'' de ``este pipeline está pronto para mapear a Amazônia'' --- a
segunda afirmação é mais forte do que os dados aqui sustentam.

O restante do artigo está organizado da seguinte forma. A
Seção~\ref{sec:dados} descreve a proveniência dos dados, o pipeline de
pré-processamento auditado, a biomassa de referência e as seis abordagens de
modelagem comparadas. A Seção~\ref{sec:resultados} apresenta os resultados: o
efeito do controle de qualidade geométrica e da correção de gap temporal, a
comparação entre arquiteturas, o diagnóstico do melhor modelo e a ablação de
dimensionalidade que motiva o argumento central do trabalho. A
Seção~\ref{sec:discussao} discute criticamente os resultados --- incluindo suas
limitações de validação --- e os posiciona frente à literatura e ao produto
orbital GEDI da NASA.

% !TEX root = ../main.tex
\section{Dados e Métodos}
\label{sec:dados}

\subsection{Proveniência dos dados}

Duas fontes brutas independentes alimentam o pipeline:

\begin{itemize}
  \item \textbf{LiDAR aerotransportado} --- nuvens de pontos \texttt{.laz} do
    programa \textit{NASA Carbon Monitoring System (CMS) Brazil}, distribuídas
    pelo ORNL DAAC. O inventário completo de tiles disponibilizado pela missão
    soma \textbf{3.152 tiles}, cobrindo múltiplas campanhas de voo (diferentes
    anos e regiões) sobre a Amazônia brasileira.
  \item \textbf{Inventário florestal de campo} --- arquivos \texttt{.csv} com
    medições árvore a árvore (diâmetro à altura do peito, altura, espécie) e
    \texttt{.kmz} com a geometria das parcelas, cobrindo \textbf{28 sites} de
    inventário distribuídos pela Amazônia.
\end{itemize}

As duas fontes são coletadas, processadas e cruzadas de forma independente uma da
outra até a etapa de interseção espacial --- nenhuma delas é derivada da outra.

\subsection{Pipeline de pré-processamento}

O cruzamento entre as duas fontes segue seis etapas, ilustradas na
Figura~\ref{fig:pipeline}.

\begin{figure}[H]
  \centering
  % PLACEHOLDER — substituir pela figura do pipeline:
  % \includegraphics[width=0.85\textwidth]{figura1_pipeline.png}
  \figplaceholder{[PLACEHOLDER] Figura 1 --- Pipeline de pré-processamento (seis
  etapas): extração de geometrias, interseção espacial parcela $\times$ tile,
  recorte do LiDAR por parcela, normalização do inventário, cálculo de biomassa e
  pareamento final.}
  \caption{Pipeline de pré-processamento que cruza o LiDAR aerotransportado com o
  inventário florestal de campo.}
  \label{fig:pipeline}
  \fonte{elaborado pelo autor (2026)}
\end{figure}

A definição de \textbf{parcela} (\texttt{plot\_id}) merece nota metodológica:
auditorias sobre os KMLs de inventário revelaram parcelas com o campo
\texttt{Name} não único --- polígonos distintos (quadras) compartilhando o mesmo
nome ---, o que levava o pipeline original a dissolver quadras diferentes numa só,
colapsando contagens e pareando AGB de campo com o clip de LiDAR errado. A
correção (documentada \textit{in extenso} no apêndice de qualidade de dados)
define o \texttt{plot\_id} por uma regra hierárquica: nome único $\rightarrow$ usa
o nome; nome repetido $\rightarrow$ sufixo sequencial (\texttt{T20\_1},
\texttt{T20\_2}, \dots); nome não informativo $\rightarrow$ índice sequencial.
Essa correção, mais a remoção de feições geometricamente duplicadas e a
recuperação de parcelas indevidamente descartadas por bugs de leitura de
coordenadas, elevou o conjunto de parcelas treináveis de 383 para as 493 finais
--- um ganho de 29\% no tamanho da amostra de treino, obtido inteiramente por
correção de bugs, sem novos dados brutos.

O critério de qualificação espacial usa \textbf{$\geq$ 99,9\% da área da parcela
dentro da cobertura contígua de uma campanha} (tiles vizinhos do mesmo voo,
unidos; subcampanhas do mesmo ano tratadas como uma só), em vez de um teste
\texttt{within} estrito --- o que evitava que parcelas 100\% cobertas, mas tocando
a borda da união de tiles, fossem indevidamente descartadas.

\subsection{Números do dataset}

\begin{table}[H]
  \centering
  \caption{Números do conjunto de dados ao longo do pipeline.}
  \label{tab:dataset}
  \begin{tabular}{lr}
    \toprule
    \textbf{Métrica} & \textbf{Valor} \\
    \midrule
    Tiles no inventário LiDAR da NASA & 3.152 \\
    Sites de inventário (total) & 28 \\
    Tiles LiDAR distintos efetivamente usados & 242 \\
    Parcelas no inventário (uma por quadra, pós-deduplicação) & 581 \\
    Pares parcela $\times$ tile & 1.287 \\
    Parcelas com cobertura LiDAR ($\geq$ 99,9\% da área) & 553 \\
    Parcelas com AGB de campo calculável (nível quadra) & 509 \\
    \textbf{Parcelas treináveis --- AGB $\times$ LiDAR pareados} & \textbf{493} \\
    Parcelas treináveis após remoção de outliers geométricos & 325 \\
    \bottomrule
  \end{tabular}
  \fonte{auditoria de dados do projeto (última varredura
  registrada), reproduzida integralmente na página de Apêndice --- Qualidade de
  Dados}
\end{table}

\subsection{Biomassa de referência (alvo)}
\label{sec:biomassa}

A AGB de campo é calculada árvore a árvore e agregada por parcela, sob três
fórmulas alométricas candidatas que diferem na especificidade da densidade de
madeira ($\rho$) e no tratamento de árvores mortas e palmeiras:

\begin{table}[H]
  \centering
  \caption{Fórmulas alométricas candidatas para a biomassa de referência.}
  \label{tab:formulas}
  \begin{tabular}{p{3.2cm}p{6.2cm}p{4.6cm}}
    \toprule
    \textbf{Fórmula} & \textbf{Equações usadas} & \textbf{Densidade de madeira ($\rho$)} \\
    \midrule
    \textbf{M1 --- Chave uniforme} & Chave et al. (2014) para todas as árvores &
      $\rho$ = 0,6 g/cm\textsuperscript{3} (média global) \\
    \textbf{M2 --- Chave por tipo} & Chave (2014, vivas) $\cdot$ Chambers (2000,
      mortas) $\cdot$ Goodman (2013, palmeiras) & $\rho$ = 0,6 g/cm\textsuperscript{3}
      (média global) \\
    \textbf{M3 --- Chave por tipo + $\rho$ específica} & Mesmas equações de M2 &
      $\rho$ por espécie (GWDD Zanne/Chave; \textit{fallback} gênero
      $\rightarrow$ site $\rightarrow$ 0,6) \\
    \bottomrule
  \end{tabular}
  \fonte{elaborado pelo autor (2026), com base em Chave et al. (2014), Chambers
  et al. (2000) e Goodman et al. (2013)}
\end{table}

Equações de carbono individual (fração de carbono $f_C = 0,5$):

\begin{align}
  \text{IAGC}_{\text{viva}}     &= 0{,}0673 \times 0{,}5 \times
     \left(\rho \cdot \text{DBH}^2 \cdot H\right)^{0{,}976}
     && \text{--- Chave et al. (2014), Eq.~7} \\
  \text{IAGC}_{\text{morta}}    &= 0{,}1007 \times 0{,}5 \times 0{,}40 \times
     \text{DBH}^2 \cdot H^{0{,}818}
     && \text{--- Chambers et al. (2000)} \\
  \text{IAGC}_{\text{palmeira}} &= 0{,}03781 \times 0{,}5 \times
     \text{DBH}^{2{,}7483}
     && \text{--- Goodman et al. (2013)}
\end{align}

com $\text{AGB} = 2 \times \text{AGC}$ (unidade final: Mg/ha). \textbf{Este
trabalho usa M1 como alvo de treino para todos os modelos.} Essa escolha --- e o
fato de não haver, nesta versão, uma análise de sensibilidade dos resultados a
M2/M3 --- é revisitada criticamente na Seção~\ref{sec:allom-sens}.

Adicionalmente, corrige-se o \textbf{descompasso temporal} entre o ano do voo
LiDAR e o ano da medição de campo: o DBH é ``crescido'' até o ano do voo pelo
incremento anual observado (por remedição, quando disponível; senão pela mediana
do site; senão um valor \textit{default}), e altura e biomassa são re-derivadas a
partir do DBH ajustado. O efeito médio é pequeno ($\sim$1,3\% na biomassa média),
mas chega a 7--12\% nos sites com gap de 2--3 anos entre voo e medição --- não
desprezível para parcelas individuais, mesmo que pequeno em média.

\subsection{Abordagens de modelagem comparadas}

Seis representações do sinal LiDAR foram comparadas, da mais compacta à mais rica
em graus de liberdade --- todas treinadas sobre o \textbf{mesmo conjunto de
parcelas}, com o mesmo protocolo de validação cruzada, para que a diferença de
desempenho seja atribuível à representação e ao modelo, não ao dado de treino:

\begin{table}[H]
  \centering
  \caption{Abordagens de modelagem comparadas.}
  \label{tab:modelos}
  \small
  \begin{tabular}{clp{5.2cm}p{3.2cm}}
    \toprule
    \textbf{\#} & \textbf{Abordagem} & \textbf{Entrada} & \textbf{Modelo} \\
    \midrule
    1 & GBR --- alturas ordenadas & 1024 alturas de dossel amostradas e ordenadas
      + log(n pontos) --- 1025 atributos & Gradient Boosting \\
    2 & GBR --- $K$ quantis & $K$ quantis igualmente espaçados da altura +
      log(n pontos) --- $K{+}1$ atributos, $K \in \{8,16,32,64,128,256\}$ &
      Gradient Boosting \\
    3 & ABA (\textit{Area-Based Approach}) & 37 métricas estruturais --- percentis
      de altura, densidade vertical, rugosidade convolucional do CHM &
      Gradient Boosting \\
    4 & Random Forest estrutural & As mesmas 37 métricas do ABA &
      Random Forest (\textit{bagging}) \\
    5 & PointNet & Nuvem de pontos crua $(x, y, \text{altura})$, até 8192
      pontos/parcela, 2048 amostrados por passo & PointNet (MLP compartilhado +
      max-pool) \\
    6 & CNN 2D --- raster & Grade 32$\times$32, 4 canais (altura máx., altura
      média, densidade log, desvio-padrão) & CNN 2D (3 blocos conv+BN+ReLU) \\
    7 & CNN 3D --- voxel & Volume 16$\times$16$\times$16, 2 canais (ocupação,
      densidade log) & CNN 3D (3 blocos conv3d+BN+ReLU) \\
    \bottomrule
  \end{tabular}
  \fonte{elaborado pelo autor (2026)}
\end{table}

As abordagens 3 e 4 compartilham exatamente as mesmas 37 features de entrada ---
a única diferença é \textit{bagging} (Random Forest) vs. \textit{boosting}
(Gradient Boosting) ---, isolando o efeito do algoritmo de ensemble do efeito da
representação. As abordagens 5--7 recebem representações geometricamente
crescentes em riqueza (nuvem crua $\rightarrow$ grade 2D $\rightarrow$ volume 3D)
e, com isso, em número de parâmetros e propensão a sobreajuste com poucas
amostras.

\subsection{Protocolo de validação}
\label{sec:protocolo}

Todos os modelos são avaliados por \textbf{validação cruzada k-fold aleatória
(5 folds)}, com o alvo transformado em log1p. ABA e Random Forest usam os mesmos
hiperparâmetros de Gradient Boosting (\texttt{n\_estimators=300}, \texttt{max\_depth=4},
\texttt{learning\_rate=0,05}, \texttt{subsample=0,8}) para que a comparação isole
o efeito do algoritmo. Reporta-se RMSE (Mg/ha), $R^2$ e rRMSE (RMSE relativo à
média observada, \%) --- agregados sobre as previsões \textit{out-of-fold} de
todos os folds (métrica ``global''), não a média simples das métricas por fold.

Cada modelo é avaliado em \textbf{seis variantes do conjunto de treino, cada uma
testando UM tratamento isoladamente contra a mesma nuvem crua} --- não empilhadas
umas sobre as outras, para que o efeito de cada tratamento possa ser atribuído a
ele mesmo, não à combinação (Seção~\ref{sec:preproc}): \textbf{sem processamento}
(\texttt{raw}, baseline), \textbf{só projeção do solo} (\texttt{ground\_only} ---
relevo local célula-a-célula, sem filtro), \textbf{só filtro $\leq$1,3~m}
(\texttt{filter\_only} --- remoção dos retornos até a altura do peito, aplicada
sobre o mesmo deslocamento mínimo de Z do baseline cru, sem acionar a projeção de
relevo completa), \textbf{só remoção de outliers} (\texttt{outliers\_only} ---
parcelas com área $>$ 0,5~ha ou compacidade $<$ 0,6 removidas, sem nenhum
tratamento na nuvem), \textbf{só correção de gap temporal} (\texttt{gap\_only} ---
alvo realinhado à data do voo LiDAR, sem nenhum tratamento na nuvem) e
\textbf{tudo combinado} (\texttt{gap\_corrected\_noout} --- os quatro tratamentos
juntos, o pipeline completo). Variantes intermediárias empilhadas (solo+filtro,
+outliers, +gap) também foram treinadas e permanecem disponíveis nos dados, mas a
análise principal da Seção~\ref{sec:preproc} usa as seis variantes isoladas como
evidência primária, por permitirem atribuir o ganho de desempenho a um tratamento
específico em vez de a uma combinação.

\begin{quote}
\textbf{Nota de escopo (ver Seção~\ref{sec:spatial-validation}):} o k-fold aqui é
aleatório sobre parcelas, não sobre sítios --- parcelas do mesmo site de
inventário podem cair em treino e teste no mesmo fold. Isso é adequado para medir
o ajuste do modelo à distribuição de estrutura do dossel presente na amostra, mas
\textbf{infla a métrica de generalização espacial} por autocorrelação entre
parcelas vizinhas do mesmo sítio --- a leitura correta dos números desta seção é
``quanto da variância intra-amostra é explicada'', não ``quão bem isto generaliza
para floresta amazônica não amostrada''.
\end{quote}

% !TEX root = ../main.tex
\section{Resultados}
\label{sec:resultados}

\subsection{Efeito do controle de qualidade dos dados}
\label{sec:preproc}

Antes de comparar arquiteturas entre si, é preciso isolar um efeito que domina os
resultados: o pré-processamento. Empilhar tratamentos uns sobre os outros (como em
versões anteriores deste texto) não permite separar o que cada etapa
individualmente contribui --- um ganho visto na combinação ``solo + filtro +
outliers'' pode vir inteiro de um único desses três, com os outros dois sendo
irrelevantes ou até prejudiciais. Para resolver isso, cada tratamento foi testado
\textbf{isoladamente}, sempre contra a mesma nuvem crua (\texttt{raw}), sem
empilhar com os demais: projeção do solo sozinha, filtro $\leq$1,3~m sozinho
(sobre o deslocamento mínimo de Z, sem o relevo local), remoção de outliers
geométricos sozinha, e correção de gap temporal sozinha no alvo. Uma sexta
coluna, ``tudo combinado'', aplica os quatro tratamentos juntos --- o pipeline
completo. A Tabela~\ref{tab:comparacao} reporta as seis colunas para os sete
modelos.

\textbf{Remoção de outliers geométricos é, disparado, o tratamento que mais
importa} --- e o único que, sozinho, já entrega quase todo o ganho visto na
combinação completa. Nos três modelos com as duas variantes disponíveis para
comparação direta (ABA, Random Forest e CNN 2D), o outlier isolado iguala ou
\textbf{supera} a combinação ``tudo junto'': CNN 2D vai de $R^2 \approx$ 0,25
(nuvem crua) para \textbf{0,55} só removendo outliers --- mais alto que os 0,48 da
combinação completa (solo+filtro+outliers+gap). Random Forest (0,53 isolado vs.
0,51 combinado) e ABA (0,51 vs. 0,49) têm o mesmo padrão. Isso indica que solo,
filtro e gap, no pipeline combinado, não estavam contribuindo --- ou estavam
discretamente atrapalhando --- e o ganho inteiro vinha da geometria das parcelas.
O GBR de 32 quantis confirma a mesma magnitude isolada ($R^2$ = 0,51 só com
outliers, Tabela~\ref{tab:comparacao}), ainda que sem uma variante ``tudo
combinado'' treinada para comparação direta.

\textbf{Projeção do solo e filtro $\leq$1,3~m, isolados, têm efeito pequeno e
inconsistente entre arquiteturas.} Só a projeção do solo (sem filtro) move o $R^2$
em poucos centésimos na maioria dos modelos (ex.: ABA 0,320$\rightarrow$0,328,
Random Forest 0,283$\rightarrow$0,286) --- exceção é a CNN 3D (voxel), que sai de
$R^2$ negativo ($-$0,13) para positivo (0,09), sugerindo que a estrutura vertical
bruta (Z absoluto, sem corrigir relevo) é especialmente ruim pra arquiteturas que
processam o volume 3D diretamente. O filtro $\leq$1,3~m isolado (sobre o mesmo
deslocamento mínimo de Z do baseline cru, sem acionar a projeção de relevo
completa) tem efeito ainda menor: piora levemente ABA, Random Forest, o modelo de
alturas ordenadas e a PointNet ($R^2$ cai alguns centésimos em cada caso ---
PointNet 0,125$\rightarrow$0,118) e ajuda discretamente só a CNN 2D e a CNN 3D,
que parecem se beneficiar de qualquer redução no ruído vertical bruto, mesmo
pequena.

\textbf{Correção de gap temporal isolada é, na maioria dos modelos, mais fraca que
não fazer nada.} Em cinco dos seis modelos com baseline cru comparável, o gap
isolado fica abaixo da nuvem crua sem nenhum tratamento --- exceção é a CNN 3D,
onde o gap isolado ($R^2 = -0,05$) é menos ruim que a própria nuvem crua
($-0,13$), embora ainda negativo. Isso contrasta com o papel do gap na combinação
completa, onde ele aparece ``ajudando'' --- mas como visto acima, o ganho ali vem
majoritariamente da remoção de outliers que o acompanha, não do gap em si. A
correção de gap é mantida no pipeline por ser metodologicamente correta (alinhar
DBH/altura do inventário à data do voo LiDAR, evitando viés sistemático de
subestimar biomassa quando o LiDAR é mais recente que o censo), não porque eleve o
$R^2$.

Em suma: \textbf{a qualidade da curadoria geométrica das parcelas (outliers)
domina a escolha de arquitetura em magnitude de efeito} --- nenhuma combinação de
solo, filtro ou gap chega perto do ganho de simplesmente descartar as parcelas de
geometria ruim. Este é o primeiro resultado do trabalho, e não é sutil.

\begin{table}[H]
  \centering
  \caption{Comparação entre modelos e variantes de pré-processamento (RMSE em
  Mg/ha; rRMSE em \%; $R^2$ adimensional; maior $R^2$ e menor RMSE/rRMSE são
  melhores). Colunas ``Só X'' são tratamentos isolados contra a mesma nuvem crua;
  ``Tudo combinado'' empilha os quatro juntos.}
  \label{tab:comparacao}
  % PLACEHOLDER — os valores por célula são gerados pelos scripts de treino
  % (data/processed/06_model/cv_results_*.csv) e exibidos no dashboard.
  % Exporte a tabela final e cole aqui no formato booktabs, ex.:
  % \begin{tabular}{lrrrrrr}
  %   \toprule
  %   Modelo & Nuvem crua & Só solo & Só filtro & Só outliers & Só gap & Tudo combinado \\
  %   \midrule
  %   ... & ... \\
  %   \bottomrule
  % \end{tabular}
  \figplaceholder{[PLACEHOLDER] Tabela 1 --- preencher com os valores de
  RMSE/$R^2$/rRMSE por modelo (7 linhas) $\times$ variante (6 colunas), exportados
  de \texttt{data/processed/06\_model/cv\_results\_*.csv}.}
  \fonte{elaborado pelo autor (2026), a partir das validações cruzadas do
  pipeline de treino}
\end{table}

\subsection{Comparação entre arquiteturas, na melhor configuração}
\label{sec:arch}

Restringindo a comparação à melhor variante de CADA modelo individualmente
($n = 325$) --- e não assumindo de antemão qual tratamento essa variante usa, já
que a Seção~\ref{sec:preproc} mostrou que cada tratamento isolado tem efeito muito
diferente ---, a variante líder é \textbf{``só outliers'' (isolada, sem
solo/filtro/gap) para cinco dos seis modelos}; CNN 3D e PointNet são as exceções,
precisando do pipeline combinado completo para atingir seu melhor resultado. A
ordem que emerge é: \textbf{CNN 2D sobre rasters de altura ($R^2$ = 0,547) $>$
Random Forest (0,532) $>$ GBR sobre a distribuição completa de 1024 alturas
(0,526) $>$ ABA (0,514) $\approx$ GBR sobre 32 quantis de altura (0,510) $>$
CNN 3D sobre voxels (0,455, pipeline combinado) $\gg$ PointNet sobre a nuvem crua
(0,212, pipeline combinado, único caso em que a variante gap-corrigida vence).}
(Esta ordem usa a fórmula alométrica de referência M1; a
Seção~\ref{sec:allom-sens} mostra que ela não é robusta à troca para M3 --- o
Random Forest assume a liderança, embora essa análise de sensibilidade ainda não
tenha sido refeita sobre as variantes isoladas identificadas aqui, ver ressalva ao
final da Seção~\ref{sec:allom-sens}.)

Essa ordem corrige uma leitura anterior deste texto, que agrupava Random Forest e
CNN 3D como equivalentes: usando a variante correta por modelo, a CNN 3D fica
\textbf{abaixo} de todos os modelos estruturais/estatísticos (ABA, GBR em ambas as
representações), não empatada com o Random Forest. A CNN 2D ainda vence por uma
margem modesta sobre Random Forest (diferença de $R^2 \approx$ 0,015 --- dentro da
variabilidade esperada de um k-fold de 5 dobras; a Seção~\ref{sec:confounders}
reporta um teste pareado formal entre CNN 2D, Random Forest e CNN 3D, cada um na
sua própria melhor variante, que não rejeita a hipótese de diferença nula entre as
três), o que sugere que a estrutura \textbf{horizontal} do dossel capturada pelo
raster (clareiras, agrupamento de copas --- informação que o ABA e o GBR de
alturas descartam por trabalhar só com a distribuição vertical agregada) carrega
sinal real, mas marginal frente à estrutura vertical já capturada pelas abordagens
mais simples.

O resultado mais informativo da tabela, porém, é o \textbf{fracasso relativo da
PointNet}: mesmo na sua melhor configuração, $R^2 \approx$ 0,21 --- pior que
\textit{todos} os outros cinco modelos, incluindo o GBR mais simples da
comparação. A PointNet recebe a representação mais rica em princípio (a nuvem de
pontos crua, sem nenhuma agregação), mas é também a que mais depende de volume de
dados para aprender uma função de agregação útil via max-pooling --- e 325
parcelas de treino são, evidentemente, poucas para isso.

\begin{figure}[H]
  \centering
  % PLACEHOLDER — substituir pelo scatter predito × observado:
  % \includegraphics[width=0.7\textwidth]{figura2_pred_obs.png}
  \figplaceholder{[PLACEHOLDER] Figura 2 --- Predito $\times$ observado, melhor
  modelo (CNN 2D, só outliers, isolado).}
  \caption{Predito $\times$ observado, melhor modelo (CNN 2D, só outliers,
  isolado). A dispersão em torno da diagonal 1:1 mostra o padrão típico de um
  modelo com $R^2$ moderado: tendência central capturada, mas variância residual
  substancial, sobretudo nas parcelas de AGB mais alta --- o regime onde erros
  absolutos (Mg/ha) pesam mais para qualquer aplicação de contabilidade de
  carbono.}
  \label{fig:pred-obs}
  \fonte{elaborado pelo autor (2026)}
\end{figure}

\subsection{Ablação de dimensionalidade: o achado central}
\label{sec:ablation}

A comparação mais direta com a pergunta de pesquisa deste trabalho não é entre
arquiteturas, mas entre \textbf{quantidade de informação de altura} fornecida ao
mesmo algoritmo (Gradient Boosting). A Figura~\ref{fig:ablation} mostra o $R^2$ de
validação cruzada em função do número de quantis de altura usados como atributo,
de 8 a 256, com o vetor completo de 1024 alturas ordenadas como referência.

O padrão contraria a intuição de que ``mais informação bruta é sempre melhor'':
\textbf{o vetor completo de 1024 alturas ($R^2 \approx$ 0,39--0,42) é superado já
pela representação mais grosseira da ablação, 8 quantis ($R^2 \approx$ 0,43)} ---
menos de 1\% da dimensionalidade original captura mais sinal preditivo, não menos.
O ganho de 8 para 16 quantis é real e consistente ($R^2$ 0,426 $\rightarrow$
0,442); daí em diante, até 64 quantis, o desempenho \textbf{satura} ($R^2$ oscila
entre 0,441 e 0,444, dentro do ruído do k-fold); só volta a subir de forma
perceptível em 256 quantis ($R^2 \approx$ 0,455) --- um ganho pequeno pelo custo de
16$\times$ mais atributos que o ponto de saturação.

\begin{figure}[H]
  \centering
  % PLACEHOLDER — substituir pela curva de ablação:
  % \includegraphics[width=0.7\textwidth]{figura3_ablation.png}
  \figplaceholder{[PLACEHOLDER] Figura 3 --- $R^2$ de validação cruzada vs. número
  de quantis de altura (escala log em x).}
  \caption{$R^2$ de validação cruzada vs. número de quantis de altura (escala log
  em $x$).}
  \label{fig:ablation}
  \fonte{elaborado pelo autor (2026)}
\end{figure}

A leitura estatística é direta: com 325 parcelas de treino e um vetor de 1024
atributos altamente correlacionados entre si (alturas vizinhas de uma distribuição
ordenada), o modelo tem mais graus de liberdade que sinal identificável --- o
excesso de dimensão é absorvido como ruído de treino (sobreajuste), não como
capacidade preditiva adicional. Reduzir para algumas dezenas de quantis não é
perda de informação relevante; é remoção de ruído correlacionado que atrapalhava a
otimização.

\subsection{O que o modelo interpretável revela}
\label{sec:interpret}

A única abordagem desta comparação que expõe explicitamente \textbf{quais
propriedades da floresta} sustentam a previsão é o ABA (GBR sobre 37 métricas
estruturais nomeadas). A Figura~\ref{fig:importance} mostra a importância das 10
variáveis mais relevantes no modelo ABA na sua variante de melhor desempenho ---
outliers isolado (Seção~\ref{sec:preproc}/\ref{sec:arch}), não o pipeline
combinado usado em versões anteriores deste texto.

\begin{figure}[H]
  \centering
  % PLACEHOLDER — substituir pelo gráfico de importância de variáveis:
  % \includegraphics[width=0.7\textwidth]{figura4_importance.png}
  \figplaceholder{[PLACEHOLDER] Figura 4 --- Importância de variáveis, modelo ABA
  (top 10).}
  \caption{Importância de variáveis, modelo ABA (top 10).}
  \label{fig:importance}
  \fonte{elaborado pelo autor (2026)}
\end{figure}

A \textbf{entropia da densidade vertical} (\texttt{vfd\_entropy}) domina
isoladamente a importância agregada ($\approx$0,34 do total), mais que o dobro da
segunda colocada --- um resultado que muda a leitura anterior deste texto,
centrada nos percentis de altura. Essa métrica mede quão uniformemente distribuída
está a folhagem ao longo do perfil vertical (entropia de Shannon sobre a densidade
de pontos por faixa de altura): um dossel com muitos estratos de folhagem bem
distribuídos tem entropia alta; um dossel concentrado numa única camada (ex.:
topo denso, meio vazio) tem entropia baixa. Apenas o percentil p75 ($\approx$0,21)
mantém posição de destaque entre os percentis de altura --- os demais (p50, p25,
p10, p05, que dominavam na variante combinada) caem para a metade inferior da
tabela. As demais variáveis relevantes --- densidade de pontos (\texttt{pt\_density}),
rugosidade convolucional do dossel (\texttt{lap\_roughness}, derivada de um
Laplaciano 3$\times$3 sobre o CHM), dispersão de altura (\texttt{h\_std}) e fração
de clareiras (\texttt{gap\_fraction}) --- carregam informação complementar sobre
heterogeneidade estrutural (vertical e horizontal) que a altura típica sozinha não
captura, e ajudam a explicar por que o raster 2D (que também enxerga
heterogeneidade horizontal diretamente, sem precisar de uma métrica resumo como
\texttt{vfd\_entropy}) supera marginalmente o ABA puro.

\subsection{Posição frente à literatura e ao GEDI}
\label{sec:gedi}

O melhor resultado desta tese ($R^2$ = 0,547; rRMSE $\approx$ 41,7\%; CNN 2D,
parcela, variante só outliers isolada) situa-se \textbf{abaixo} dos trabalhos de
referência em LiDAR aéreo tropical com metodologia mais madura --- Asner \&
Mascaro (2014), $R^2$ = 0,90 sobre 804 parcelas de 1~ha em quatro países tropicais
e Havaí; Longo et al. (2016), $R^2$ = 0,70 sobre 359 parcelas amazônicas --- e
também \textbf{abaixo} de Bispo et al. (2020): $R^2$ = 0,93 / rRMSE = 13\% para o
modelo só-LiDAR ($R^2$ = 0,89 / rRMSE = 43\% no modelo multissensor). \textit{Nota:
Bispo et al. (2020) estuda o Cerrado brasileiro (bacia do Rio Vermelho), não a
Amazônia --- versões anteriores deste texto atribuíam o estudo à Amazônia e
citavam $R^2$ = 0,58 / rRMSE = 36,0\%, que não correspondem à fonte; ambos os
erros foram corrigidos após verificação adversarial direta contra o artigo
original.} Frente ao produto orbital \textbf{GEDI L4A} (\textit{footprint}
$\sim$25~m), a validação independente em savanas do sul da África (Li et al., 2024
--- não Naidoo et al., que é coautor, não autor líder; ver correção na seção de
Referências) reporta $R^2 \approx$ 0,42 --- nesse critério específico, o resultado
desta tese é \textbf{competitivo}, apesar de operar em escala de parcela (mais
exigente) em vez de \textit{footprint} de 25~m. d'Oliveira et al. (2012), com
suporte espacial menor (50 parcelas de 0,25~ha, Acre) e instrumentação mais
antiga, reporta $R^2$ = 0,63--0,72 --- o mais baixo do conjunto de referências ---
e foi mantido na lista de referências sem entrar nesta comparação direta por essas
duas diferenças de desenho.

Essas comparações exigem duas ressalvas, detalhadas na Seção~\ref{sec:lit-check}:
(i) as escalas espaciais não são diretamente comparáveis (parcela vs.
\textit{footprint} vs. grade de 1~km); (ii), mais importante, os números desta
tese vêm de k-fold aleatório, enquanto boa parte da literatura citada usa
validação cruzada por sítio ou validação independente --- um padrão de
generalização mais exigente que o empregado aqui (Seção~\ref{sec:spatial-validation}).
A verificação adversarial dos números da literatura --- conferência linha a linha
contra a fonte primária, antes pendente --- foi concluída; ver
Seção~\ref{sec:lit-check} para o detalhe por fonte.

Quanto ao LiDAR orbital da NASA --- GEDI e ICESat-2 ---, não foi encontrada, mesmo
após busca e verificação direta contra a fonte primária, nenhuma validação de
$R^2$/RMSE calculada isoladamente para a Amazônia brasileira (nem do GEDI, nem do
ICESat-2) com desenho comparável a parcela ou k-fold aleatório. O benchmark
publicado mais próximo é Duncanson et al. (2022): o modelo oficial de biomassa do
GEDI, no estrato \textit{Evergreen Broadleaf Trees} --- América do Sul (classe
dominante na Amazônia, mas estrato continental, não um recorte amazônico), reporta
$R^2$ = 0,66 e \%RMSE = 42,2\% por validação cruzada geográfica --- portanto nem
k-fold aleatório, nem especificamente amazônico, mas é o melhor proxy disponível
para ``GEDI num bioma análogo à Amazônia'', na mesma faixa de rRMSE desta tese
(41,7\%). Já o ICESat-2 (ATL08), quando validado contra LiDAR aerotransportado em
floresta tropical (Mesoamérica, Gabão e Guiana Francesa --- nunca a bacia
amazônica isolada) (Fernandez-Diaz; Velikova; Glennie, 2022), degrada fortemente frente ao desempenho em floresta boreal:
erro de altura de dossel de 5--7~m em floresta tropical contra 3,19~m em condições
boreais ideais (Liu, Gong \& Shi, 2026) --- reforçando que mesmo o LiDAR orbital
mais recente da NASA perde precisão em dossel tropical denso, um argumento a favor
da abordagem aerotransportada adotada nesta tese. O trabalho de deep learning mais
próximo publicado em 2023--2025 usando a mesma linhagem de dados (LiDAR
aerotransportado do programa NASA CMS Brasil/EBA) na Amazônia é Wagner et al.
(2025, pré-print): uma U-Net treinada com esse mesmo tipo de levantamento para
mapear altura de dossel (não AGB) a partir de imagens Planet NICFI, com
MAE = 3,68~m em validação espacialmente estratificada --- uma comparação
metodológica (arquitetura de rede) mais honesta que uma comparação direta de
métrica de erro, já que a variável-alvo difere da desta tese.

% !TEX root = ../main.tex
\section{Discussão}
\label{sec:discussao}

\subsection{O achado central: parcimônia supera profundidade}

O resultado mais robusto deste trabalho não é o $R^2$ = 0,53 da CNN 2D --- é a
constatação, convergente por duas vias independentes (ablação de dimensionalidade,
Seção~\ref{sec:ablation}; e importância de variáveis do modelo interpretável,
Seção~\ref{sec:interpret}), de que \textbf{a distribuição de altura do dossel,
resumida em poucas dezenas de números, carrega quase todo o sinal preditivo de AGB
disponível nesta amostra}. Arquiteturas de maior capacidade representacional
(PointNet sobre nuvem crua, CNN 3D sobre voxels) não convertem essa riqueza
geométrica adicional em desempenho preditivo superior --- pelo contrário, a
PointNet é consistentemente a pior de todas, e a CNN 3D só se torna competitiva
depois de um controle de qualidade de dados agressivo. Este é o enquadramento
correto da contribuição do trabalho: não ``qual arquitetura é a melhor para AGB por
LiDAR'', mas \textbf{em que ponto a complexidade do modelo deixa de ajudar e passa
a atrapalhar, dado um orçamento de poucas centenas de parcelas rotuladas} --- um
regime realista para a maior parte dos estudos de campo na Amazônia, onde
inventários de centenas de parcelas já representam um esforço logístico
considerável.

\subsection{A limitação mais importante: validação espacialmente ingênua}
\label{sec:spatial-validation}

A validação cruzada k-fold aleatória usada em todos os resultados da
Seção~\ref{sec:resultados} particiona \textbf{parcelas}, não \textbf{sítios}. É
importante separar dois mecanismos distintos pelos quais isso pode inflar o $R^2$
reportado --- eles têm força de evidência diferente, e merecem tratamento
diferente.

\textbf{(a) Pseudo-replicação dentro de uma mesma quadra/campo --- mecanismo que
não depende de heterogeneidade biológica entre sítios.} Muitas das 325 parcelas do
conjunto final não são amostras independentes de ``a floresta amazônica'': são
\textbf{subdivisões contíguas de um mesmo campo original}. Dos 325, \textbf{125
(38\%) pertencem a um grupo de $\geq$2 quadras-irmãs} derivadas do mesmo
\texttt{Name} de campo (ex.: CAU tem campos como T17--T22 cada um dividido em 4--5
quadras adjacentes; FST tem 4 campos divididos em 5 quadras cada) --- e \textbf{um
único sítio (CAU) responde por 85 das 325 parcelas (26\% da amostra inteira)}.
Quadras vizinhas de um mesmo campo compartilham, literalmente, o mesmo dossel
contínuo ou dosséis a poucas dezenas de metros um do outro: seus vetores de altura
são amostras altamente correlacionadas da mesma realização física de floresta, não
realizações independentes de ``estrutura típica da Amazônia''. Um modelo pode
ajustar o padrão numérico específico de um campo (treino) e prever bem sua
quadra-irmã (teste) sem ter aprendido nada que se transporte para um hectare
genuinamente não amostrado --- isso é verdade \textbf{mesmo sob a hipótese de que
toda a floresta amazônica fosse biologicamente homogênea}, porque o problema não é
``os sítios são diferentes entre si'', é ``as unidades de teste e treino, dentro de
um sítio, não são estatisticamente independentes''. Esse mecanismo é o mais sólido
dos dois e o que mais pesa nesta ressalva.

\textbf{(b) Confundidor de campanha/densidade de varredura --- mecanismo
condicional, cuja força depende de quanto os sítios de fato diferem entre si.} Se a
floresta amostrada for de fato aproximadamente homogênea entre sítios (mesmo bioma,
mesma faixa de biomassa típica, sem efeito sistemático de solo/perturbação/região
correlacionado com o alvo) e nenhuma feature identificar o sítio diretamente, este
segundo mecanismo é fraco --- e a crítica de que ``misturar sítios é um problema''
perde força nesse cenário, como argumentado. Mas a premissa de homogeneidade não é
demonstrada pelos próprios dados do projeto: o pipeline de alometria
(\texttt{config/allometry.yaml}) usa curvas altura-diâmetro \textbf{regionais
distintas} (Amazônia centro-oriental, ocidental, Escudo das Guianas, Escudo
Brasileiro) precisamente porque a relação estrutura$\leftrightarrow$biomassa varia
por região --- e a feature $\log(\text{n}^{\text{o}}\text{ de pontos})$ /
\texttt{pt\_density}, presente em todos os modelos, captura densidade de varredura,
que correlaciona com campanha de voo e, por acaso, com região. Este mecanismo
permanece \textbf{plausível, não testado} (Seção~\ref{sec:confounders}) --- mais
fraco que (a), mas não eliminado.

A consequência prática: mesmo que (b) se revele desprezível após a análise
estratificada por campanha (Seção~\ref{sec:confounders}, ainda pendente), o
mecanismo (a) por si só já é suficiente para que os números da
Seção~\ref{sec:resultados} devam ser lidos como ``quanto da variância é explicada
quando se permite reaproveitar informação de quadras vizinhas do mesmo campo'' ---
não ``quão bem isto generaliza para um hectare novo da Amazônia'', que é a pergunta
que a motivação do trabalho (Seção~\ref{sec:intro}) coloca. Fechar essa lacuna
exige reportar uma validação agrupada por sítio (ou, de forma ainda mais
conservadora, por campo/\texttt{Name} de origem da quadra) como métrica principal,
não como nota de rodapé --- com a ressalva prática de que os 18 sítios do conjunto
final têm tamanhos muito desiguais (de 4 a 85 parcelas), então um
\textit{leave-one-site-out} ingênuo produzirá folds de variância muito heterogênea
e deve ser reportado por sítio, não só agregado. É esperável que o $R^2$ caia sob
esse protocolo mais exigente. Isto não invalida o achado central da
Seção~\ref{sec:discussao} (a relação entre dimensionalidade e desempenho deve se
manter qualitativamente, já que ambos os mecanismos acima afetam todos os modelos
comparados de forma semelhante), mas invalida qualquer leitura dos números
absolutos de $R^2$/RMSE como medida de prontidão operacional.

\subsection{Quantificação de incerteza: intervalo conforme por parcela}
\label{sec:conformal}

A motivação declarada deste trabalho --- contabilidade de carbono, REDD+, mercados
de crédito --- é um domínio em que uma estimativa pontual de AGB sem intervalo de
predição tem utilidade limitada: qualquer aplicação de compensação ou relatório de
emissões exige saber não só o valor central, mas o \textbf{intervalo de confiança}
da estimativa por parcela ou por pixel. As métricas reportadas neste artigo
($R^2$, RMSE, rRMSE) descrevem o comportamento agregado do modelo sobre o conjunto
de teste; nenhuma delas produz, por si só, um intervalo de predição individual.

Um intervalo conforme por split (\textit{split conformal prediction}) foi
calculado sobre os resíduos out-of-fold já salvos de cada modelo (variante sem
outliers, o mesmo conjunto reportado na Seção~\ref{sec:arch}): os resíduos
absolutos de metade das parcelas (calibração, $n \approx$ 162) definem o raio do
intervalo a 90\% de confiança; a outra metade, nunca usada na calibração, mede a
cobertura empírica real. O resultado: raios de $\pm$109~Mg/ha (CNN 2D, o modelo
mais acurado) a $\pm$144~Mg/ha (PointNet, o menos acurado), com cobertura empírica
entre 86\% e 92\% --- próxima da nominal de 90\% em todos os seis modelos, o que
indica que o intervalo conforme é razoavelmente bem calibrado nesta amostra,
apesar do tamanho reduzido da calibração. Um raio de $\pm$109 a $\pm$144~Mg/ha
sobre uma AGB média de $\approx$ 170~Mg/ha (variante sem outliers) é, na prática,
um intervalo largo --- a maior parte da AGB plausível de uma parcela amazônica cai
dentro dele --- mas é a primeira quantificação honesta de quão incerta é cada
estimativa individual, ao contrário de reportar só $R^2$/RMSE agregados. Este
intervalo cobre apenas a incerteza do modelo estrutura$\rightarrow$AGB; não propaga
o erro alométrico da Seção~\ref{sec:biomassa} nem a incerteza de medição de campo
até o alvo de treino --- ambos permanecem como trabalho futuro.

\subsection{O alvo de treino é, ele mesmo, um modelo}
\label{sec:allom-sens}

A AGB de campo usada como alvo de treino não é uma medição direta: é a saída de
uma equação alométrica (Seção~\ref{sec:biomassa}) aplicada a medições de campo
(DBH, altura, espécie), e essa equação tem incerteza própria, documentada na
literatura que a origina (Chave et al., 2014). Este trabalho reporta os resultados
principais (Seção~\ref{sec:resultados}) usando \textbf{M1} (densidade de madeira
uniforme, $\rho$ = 0,6 g/cm\textsuperscript{3}) como alvo --- a escolha mais
simples entre as três fórmulas calculadas na Seção~\ref{sec:biomassa} (M2 por tipo
de árvore, M3 por densidade específica de espécie).

Uma análise de sensibilidade foi conduzida: os sete modelos foram retreinados
contra \textbf{M3} (densidade de madeira específica por espécie) sob o mesmo
protocolo de validação cruzada.

\textbf{Ressalva de leitura}: a coluna $R^2$ (M3) abaixo foi treinada quando a
variante de melhor desempenho conhecida para cada modelo ainda era o pipeline
combinado (``solo+filtro+outliers'', sem gap) --- a Seção~\ref{sec:preproc}
identificou depois, via ablação isolada, que a remoção de outliers SOZINHA (sem
solo/filtro) é a variante líder para cinco dos seis modelos, superando inclusive o
combinado em alguns casos. A comparação M1 vs. M3 abaixo ainda não foi refeita
sobre essas variantes isoladas --- a coluna $R^2$ (M1) foi atualizada para
refletir os novos melhores valores (Seção~\ref{sec:arch}), mas a coluna $R^2$ (M3)
permanece com o valor do pipeline combinado antigo; $\Delta R^2$ é, portanto, um
limite inferior pessimista da perda real sob M3 nos modelos onde a variante
isolada lidera.

\begin{table}[H]
  \centering
  \caption{Sensibilidade do desempenho à fórmula alométrica de referência
  (M1 vs. M3).}
  \label{tab:m1m3}
  \begin{tabular}{lccc}
    \toprule
    \textbf{Modelo} & \textbf{$R^2$ (M1)} &
      \textbf{$R^2$ (M3, pipeline combinado)} & \textbf{$\Delta R^2$} \\
    \midrule
    CNN 2D                 & 0,547 & 0,286 & \textbf{$-$0,261} \\
    Random Forest          & 0,532 & 0,382 & $-$0,150 \\
    GBR alturas completas  & 0,526 & 0,309 & $-$0,217 \\
    ABA                    & 0,514 & 0,360 & $-$0,154 \\
    GBR 32 quantis         & 0,510 & 0,329 & $-$0,181 \\
    CNN 3D                 & 0,455 & 0,327 & $-$0,128 \\
    PointNet               & 0,212 & 0,128 & $-$0,084 \\
    \bottomrule
  \end{tabular}
  \fonte{elaborado pelo autor (2026), a partir das validações cruzadas do
  pipeline de treino sob M1 e M3}
\end{table}

Dois achados, nenhum deles trivial. Primeiro, \textbf{todos os sete modelos perdem
poder preditivo substancial sob M3} ($\Delta R^2$ entre $-$0,08 e $-$0,26) --- a
densidade de madeira específica por espécie não é visível na estrutura do dossel
medida pelo LiDAR (nenhuma das representações usadas aqui recebe identidade
taxonômica como insumo), então parte da variância que M3 introduz no alvo é, para
qualquer um destes modelos, ruído irredutível. Segundo, e mais importante para a
validade do resultado central da Seção~\ref{sec:arch}: \textbf{a escolha da fórmula
muda a ordem dos modelos}. Sob M1, a CNN 2D lidera; sob M3 (na comparação
disponível, pipeline combinado), ela cai para o penúltimo lugar (só acima da
PointNet), e o Random Forest assume a liderança --- a maior queda absoluta de toda
a comparação ($\Delta R^2 \approx -0,26$) é justamente do modelo antes apontado
como o melhor. Isso não invalida o argumento de que arquiteturas parcimoniosas
(ABA, Random Forest, GBR) superam a PointNet --- essa conclusão sobrevive intacta
sob as duas fórmulas ---, mas enfraquece a afirmação mais específica de que a
CNN 2D é a melhor arquitetura desta comparação: esse resultado é sensível à
escolha, metodologicamente defensável mas arbitrária entre três candidatas
igualmente válidas, de qual fórmula alométrica usar como alvo. \textbf{Trabalho
pendente}: retreinar a comparação M3 sobre a variante isolada de outliers (líder
sob M1 para cinco dos seis modelos) para confirmar se a inversão de ordem persiste
nessa variante também.

\subsection{Confundidores de campanha e densidade de varredura}
\label{sec:confounders}

As nuvens de pontos usadas neste trabalho vêm de campanhas de voo distintas ---
anos diferentes, sensores possivelmente distintos, densidades de varredura
variáveis --- e uma das features usadas por praticamente todas as abordagens
($\log(\text{n}^{\text{o}}\text{ de pontos})$ ou \texttt{pt\_density}) captura
diretamente essa densidade de aquisição, não uma propriedade da floresta. Isso
abre a possibilidade de que parte do $R^2$ reportado reflita o modelo aprendendo a
\textbf{identificar a campanha} (que pode estar correlacionada com a região e, por
acaso, com a biomassa típica daquela região) em vez de aprender a relação
estrutura $\leftrightarrow$ biomassa que se pretende generalizar.

Duas análises, conduzidas sobre o ABA na sua variante de melhor desempenho
(outliers isolado --- Seção~\ref{sec:preproc}/\ref{sec:arch}), testam essa
hipótese. \textbf{$R^2$ estratificado por ano de campanha} varia bastante --- de
0,12 (2014, $n$=107) a 0,39 (2013, $n$=120), passando por um valor negativo em
2016 ($n$=9, amostra pequena demais para ser conclusiva --- ver ressalva da
Seção~\ref{sec:conformal} sobre fold sizes desbalanceados). Essa variação é
consistente com a hipótese do confundidor, mas não a confirma isoladamente:
campanhas diferentes também cobrem florestas estruturalmente diferentes, então
parte da variação pode ser heterogeneidade biológica real, não viés de densidade
de varredura. Um \textbf{modelo de controle sem as duas features que capturam
densidade de varredura diretamente} (\texttt{pt\_density}, \texttt{ptdens\_cv})
perde pouco desempenho ($R^2$ 0,514 $\rightarrow$ 0,496, $\Delta R^2 \approx$ 0,02)
--- a maior parte do sinal do modelo completo sobrevive sem essas features, o que
pesa contra (mas não elimina) a hipótese de que o modelo está majoritariamente
``reconhecendo a campanha'' pela densidade de pontos. As demais features
(percentis de altura, deciles de densidade por faixa de altura \texttt{D0}--\texttt{D9},
métricas convolucionais do CHM) também podem carregar correlação residual com a
campanha de forma mais sutil, que esta análise não isola.

\subsection{Qualidade do rótulo e heterogeneidade de área}

A auditoria de dados (Apêndice --- Qualidade de Dados) documenta parcelas de área
muito pequena ($\sim$0,02~ha) gerando AGB por hectare implausivelmente alta ---
acima de 1500~Mg/ha, mais que o dobro do teto fisicamente plausível para floresta
tropical densa ---, um artefato estatístico de dividir uma contagem pequena de
árvores por uma área minúscula. O critério de remoção de outliers usado nesta
versão (área $>$ 0,5~ha OU compacidade $<$ 0,6) mitiga a maior parte do problema: a
área das 305 parcelas remanescentes vai de 0,077~ha a 0,305~ha, não mais os
0,02~ha a 1~ha do conjunto completo. Mesmo nessa faixa mais estreita, uma AGB por
hectare estimada sobre poucas árvores tem variância amostral estruturalmente maior
que a mesma métrica sobre uma parcela de 1~ha --- o que significa que as métricas
agregadas de RMSE/$R^2$ da Seção~\ref{sec:resultados} não são igualmente
confiáveis para todas as parcelas do conjunto, mas isto não é reportado de forma
estratificada.

\subsection{Números da literatura: verificação concluída (três erros corrigidos)}
\label{sec:lit-check}

A comparação da Seção~\ref{sec:gedi} usava, em versões anteriores deste texto,
valores extraídos de fontes primárias por pesquisa automatizada sem verificação
adversarial linha a linha. Essa verificação foi concluída e encontrou três erros,
já corrigidos no texto atual: (i) Bispo et al. (2020) estava atribuído à Amazônia
com $R^2$ = 0,58 / rRMSE = 36,0\% --- o estudo é sobre o Cerrado brasileiro e os
valores reais são $R^2$ = 0,93 / rRMSE = 13\% (modelo só-LiDAR); (ii) o DOI citado
para Asner \& Mascaro (2014) apontava para um artigo diferente (albedo boreal,
mesma revista/volume); (iii) o número do GEDI L4A em savana africana
($R^2 \approx$ 0,42, correto) estava atribuído a ``Naidoo et al.'' --- o autor
líder correto é Li et al. (2024), com Naidoo como coautor. Longo et al. (2016) e
Oehmcke et al. (2024) foram confirmados sem ressalva. Esta correção ilustra
exatamente o risco que a versão anterior deste parágrafo já havia antecipado: um
número mal transcrito --- ou, neste caso, três --- compromete a credibilidade de
toda a comparação até ser verificado.

\subsection{Síntese e trabalho futuro}

Em ordem de prioridade para que este trabalho evolua de exploração metodológica
para resultado defensável de dissertação:

\begin{enumerate}
  \item \textbf{Validação LOSO} (Seção~\ref{sec:spatial-validation}) como métrica
    principal, k-fold aleatório rebaixado a referência otimista secundária.
  \item \textbf{Intervalos de predição} (Seção~\ref{sec:conformal}) --- bootstrap
    ou conformal, com cobertura empírica reportada.
  \item \textbf{Análise de sensibilidade alométrica} (Seção~\ref{sec:allom-sens})
    --- resultados da Seção~\ref{sec:resultados} repetidos sob M2 e M3.
  \item \sout{\textbf{Verificação adversarial da literatura}
    (Seção~\ref{sec:lit-check}) antes de qualquer citação em versão final.}
    Concluída --- três erros corrigidos (ver Seção~\ref{sec:lit-check}).
  \item \textbf{Estratificação por campanha e por tamanho de parcela}
    (Seções~\ref{sec:confounders} e seguinte) para descartar confundidores e
    reportar honestamente onde o modelo é, e não é, confiável.
  \item Avaliar formalmente o \textit{ensemble} (média simples de raster + RF +
    voxel, que a exploração preliminar do projeto situa acima de qualquer modelo
    isolado) sob o mesmo protocolo LOSO --- um ganho medido sob k-fold aleatório
    pode não sobreviver ao protocolo mais exigente.
\end{enumerate}

Nenhum desses pontos invalida o achado central da Seção~\ref{sec:discussao}. Pelo
contrário: se a relação entre parcimônia e desempenho se mantiver sob validação
LOSO --- o teste mais rigoroso disponível com os dados atuais ---, ela se torna uma
contribuição metodológica transferível para outros contextos de estimativa de
biomassa por LiDAR com amostras de campo reduzidas, para além do caso amazônico
específico deste trabalho.


\makeatletter
\renewcommand{\@biblabel}[1]{}
\makeatother
\begin{thebibliography}{99}
\setlength{\itemsep}{1em}
\setlength{\parindent}{0pt}
\renewcommand{\bibname}{Referências}

\bibitem{asner2014} ASNER, G. P.; MASCARO, J. Mapping tropical forest carbon:
calibrating plot estimates to a simple LiDAR metric. \textit{Remote Sensing of
Environment}, v. 140, p. 614-624, 2014. Disponível em:
\url{https://doi.org/10.1016/j.rse.2013.09.023}. Acesso em: 3 nov. 2026.

\bibitem{bispo2020} BISPO, P. C. et al. Woody aboveground biomass mapping of the
Brazilian savanna with a multi-sensor and machine learning approach.
\textit{Remote Sensing of Environment}, v. 232, p. 111453, 2020. Disponível em:
\url{https://doi.org/10.1016/j.rse.2019.111453}. Acesso em: 3 nov. 2026.

\bibitem{chambers2000} CHAMBERS, J. Q.; HIGUCHI, N.; SCHIMEL, J. P.; FERREIRA,
L. V.; MELACK, J. M. Decomposition and carbon cycling of dead trees in tropical
forests of the central Amazon. \textit{Oecologia}, v. 122, p. 380-388, 2000.

\bibitem{chave2014} CHAVE, J.; RÉJOU-MÉCHAIN, M.; BÚRQUEZ, A.; CHIDUMAYO, E.;
COLGAN, M. S.; DELITTI, W. B. C.; DUQUE, A.; EID, T.; FEARNSIDE, P. M.; GOODMAN,
R. C.; HENRY, M.; MARTÍNEZ-YRÍZAR, A.; MUGASHA, W. A.; MULLER-LANDAU, H. C.;
MENCUCCINI, M.; NELSON, B. W.; NGOMANDA, A.; NOGUEIRA, E. M.; ORTIZ-MALAVASSI,
E.; PÉLISSIER, R.; PLOTON, P.; RYAN, C. M.; SALDARRIAGA, J. G.; VIEILLEDENT, G.
Improved allometric models to estimate the aboveground biomass of tropical
trees. \textit{Global Change Biology}, v. 20, n. 10, p. 3177-3190, 2014.

\bibitem{doliveira2012} D'OLIVEIRA, M. V. N.; REUTEBUCH, S. E.; McGAUGHEY,
R. J.; ANDERSEN, H.-E. Estimating forest biomass and identifying low-intensity
logging areas using airborne scanning lidar in Antimary State Forest, Acre
State, Western Brazilian Amazon. \textit{Remote Sensing of Environment},
v. 124, p. 479-491, 2012. Disponível em:
\url{https://doi.org/10.1016/j.rse.2012.05.014}. Acesso em: 3 nov. 2026.

\bibitem{duncanson2022} DUNCANSON, L. et al. Aboveground biomass density
models for NASA's Global Ecosystem Dynamics Investigation (GEDI) lidar
mission. \textit{Remote Sensing of Environment}, v. 270, p. 112845, 2022.
Disponível em: \url{https://doi.org/10.1016/j.rse.2021.112845}. Acesso em:
3 nov. 2026.

\bibitem{fao2020} FAO. \textit{Global Forest Resources Assessment 2020: main
report}. Roma: FAO, 2020. Disponível em:
\url{https://doi.org/10.4060/ca9825en}. Acesso em: 3 nov. 2026.

\bibitem{fernandezdiaz2022} FERNANDEZ-DIAZ, J. C.; VELIKOVA, M.; GLENNIE, C. L.
Validation of ICESat-2 ATL08 terrain and canopy height retrievals in tropical
Mesoamerican forests. \textit{IEEE Journal of Selected Topics in Applied Earth
Observations and Remote Sensing}, v. 15, 2022. Disponível em:
\url{https://doi.org/10.1109/JSTARS.2022.3163208}. Acesso em: 3 nov. 2026.

\bibitem{gediL4A} GEDI L4A. \textit{Footprint Level 4A Aboveground Biomass
Density}, Version 2.1. Oak Ridge, TN: ORNL DAAC, 2022. Disponível em:
\url{https://daac.ornl.gov/GEDI/guides/GEDI_L4A_AGB_Density_V2_1.html}. Acesso
em: 3 nov. 2026.

\bibitem{gediL4B} GEDI L4B. \textit{Gridded Aboveground Biomass Density},
Version 2.1. Oak Ridge, TN: ORNL DAAC, 2022. Disponível em:
\url{https://daac.ornl.gov/GEDI/guides/GEDI_L4B_Gridded_Biomass_V2_1.html}.
Acesso em: 3 nov. 2026.

\bibitem{goodman2013} GOODMAN, R. C.; PHILLIPS, O. L.; DEL CASTILLO TORRES, D.;
FREITAS, L.; CORTESE, S. T.; MONTEAGUDO, A.; BAKER, T. R. Amazon palm biomass
and allometry. \textit{Forest Ecology and Management}, v. 310, p. 994-1004,
2013.

\bibitem{ipcc2019} IPCC. \textit{2019 Refinement to the 2006 IPCC Guidelines
for National Greenhouse Gas Inventories}. Genebra: IPCC, 2019. Disponível em:
\url{https://www.ipcc-nggip.iges.or.jp/public/2019rf/index.html}. Acesso em:
3 nov. 2026.

\bibitem{ipcc2021} IPCC. \textit{Climate Change 2021: The Physical Science
Basis. Contribution of Working Group I to the Sixth Assessment Report of the
Intergovernmental Panel on Climate Change}. Cambridge: Cambridge University
Press, 2021. Disponível em: \url{https://www.ipcc.ch/report/ar6/wg1/}. Acesso
em: 3 nov. 2026.

\bibitem{li2024} LI, X.; WESSELS, K.; ARMSTON, J.; DUNCANSON, L.; URBAZAEV, M.;
NAIDOO, L.; MATHIEU, R.; MAIN, R. Evaluation of GEDI footprint level biomass
models in Southern African savannas using airborne LiDAR and field
measurements. \textit{Science of Remote Sensing}, v. 10, p. 100161, 2024.
\textit{Nota}: citado em versões anteriores deste texto como ``Naidoo et al.''
--- Naidoo é coautor, não autor líder; corrigido após verificação adversarial
contra a fonte primária.

\bibitem{liu2026} LIU, A.; GONG, P.; SHI, Y. Evaluating and calibrating
ICESat-2 canopy height: airborne validation and machine learning enhancement
across boreal and tropical forests. \textit{Forests}, v. 17, n. 2, p. 185,
2026. Disponível em: \url{https://doi.org/10.3390/f17020185}. Acesso em:
3 nov. 2026.

\bibitem{longo2016} LONGO, M. et al. Aboveground biomass variability across
intact and degraded forests in the Brazilian Amazon. \textit{Global
Biogeochemical Cycles}, v. 30, n. 11, p. 1639-1660, 2016. Disponível em:
\url{https://doi.org/10.1002/2016GB005465}. Acesso em: 3 nov. 2026.

\bibitem{oehmcke2024} OEHMCKE, S.; LI, L.; TREPEKLI, K.; REVENGA, J. C.;
NORD-LARSEN, T.; GIESEKE, F.; IGEL, C. Deep point cloud regression for
above-ground forest biomass estimation from airborne LiDAR. \textit{Remote
Sensing of Environment}, v. 302, p. 113968, 2024. Disponível em:
\url{https://doi.org/10.1016/j.rse.2023.113968}. Acesso em: 3 nov. 2026.

\bibitem{pan2011} PAN, Y. et al. A large and persistent carbon sink in the
world's forests. \textit{Science}, v. 333, n. 6045, p. 988-993, 2011.
Disponível em: \url{https://doi.org/10.1126/science.1201609}. Acesso em:
3 nov. 2026.

\bibitem{saatchi2011} SAATCHI, S. S. et al. Benchmark map of forest carbon
stocks in tropical regions across three continents. \textit{Proceedings of
the National Academy of Sciences}, v. 108, n. 24, p. 9899-9904, 2011.
Disponível em: \url{https://doi.org/10.1073/pnas.1019576108}. Acesso em:
3 nov. 2026.

\bibitem{unfccc2015} UNFCCC. \textit{Paris Agreement}. Paris: United Nations,
2015. Disponível em:
\url{https://unfccc.int/process-and-meetings/the-paris-agreement}. Acesso em:
3 nov. 2026.

\bibitem{wagner2025} WAGNER, F. H.; DALAGNOL, R.; CARTER, G.; KELLER, M.;
SAATCHI, S. \textit{High resolution tree height mapping of the Amazon forest
using Planet NICFI images and LiDAR-informed U-Net model}. arXiv preprint,
arXiv:2501.10600, 2025. \textit{Nota}: pré-print --- não confirmada publicação
revisada por pares sob este título; usa LiDAR aerotransportado do mesmo
programa NASA CMS Brasil/EBA desta tese, mas estima altura de dossel, não AGB.

\end{thebibliography}


\vspace{1em}
\noindent\textbf{Fontes de dados.}
NASA Carbon Monitoring System (CMS) Brazil --- LiDAR aerotransportado, ORNL DAAC;
inventários florestais de campo integrados ao mesmo programa (ORNL DAAC), medições
árvore a árvore (DBH, altura, espécie) e geometria de parcela (.kmz/.kml).

\end{document}
